\section{Introduction}\label{sec:intro}
Perception models for autonomous driving need large, varied and labelled LiDAR data, which is costly to record and annotate. Synthetic data offers control over scenes and automatic labels, but synthetic scans differ from real measurements, and this domain gap can limit transfer~\cite{hurlPreciseSyntheticImage2019,manivasagamLiDARsimRealisticLiDAR2020}. Our literature review found that PreSIL, LiDARsim and SynLiDAR show useful transfer in particular settings~\cite{hurlPreciseSyntheticImage2019,manivasagamLiDARsimRealisticLiDAR2020,xiaoTransferLearningSynthetic2022}. It also found that the most direct evidence of training utility is the performance on held-out real data of a model trained on the synthetic data.

That evidence is expensive: a model has to be trained and evaluated for every candidate dataset or sensor setting. Cheaper data-level diagnostics are therefore widely used to judge generated scans~\cite{zyrianovLearningGenerateRealistic2022}. Three of them are central here. Each compares a set of synthetic scans with a set of real scans and returns one number, where a lower value means the two sets are more alike:
\begin{itemize}
    \item The \emph{Fr\'echet range distance (FRD)} passes real and synthetic scans through a pretrained segmentation network (RangeNet++) and compares the average and the spread of the features it extracts. It asks whether the network ``sees'' the two sets in the same way.
    \item The \emph{maximum mean discrepancy (MMD)} compares the two sets using bird's-eye-view (BEV) maps, which show how many points fall in each cell of a top-down grid. It is small when the average real and synthetic maps are hard to tell apart.
    \item The \emph{Jensen--Shannon divergence (JSD)} treats the same BEV maps as probability distributions over the grid and measures how much they differ. It is 0 when the distributions are identical.
\end{itemize}
Performance on real data is measured with the mean intersection over union (mIoU): for each class, the overlap between predicted and true labels divided by their union, averaged over classes. Unlike the three diagnostics, mIoU can also be reported per class and per distance range.

This proposal studies how sensitive the diagnostics are to controlled sensor degradations that reduce real-data performance. Section~\ref{sec:problem} states the problem, Section~\ref{sec:questions} the research questions, Section~\ref{sec:hypotheses} the hypotheses, and Sections~\ref{sec:method}--\ref{sec:conditions} the method, plan and conditions.